\section*{Chapitre \ref{ch:pl:trade-capture-and-booking} --- Saisie des transactions et enregistrement}

\begin{solution}{exo:pl:trade-capture-and-booking:1}
Dix lots, parts exactes de $3\tfrac13$ lots chacune. Plus grand reste~: 400, 300 et 300 (l'égalité étant départagée par les noms des fonds). Lots irréguliers attribués aux
plus grands fonds~: même résultat, puisque les poids sont égaux. Arrondi de chaque fonds~: 300 chacun, soit 900 actions, ce qui laisse 100 actions non attribuées.
\end{solution}

\begin{solution}{exo:pl:trade-capture-and-booking:2}
Il copie chaque transaction chaque soir, de sorte qu'il ne manque que ce qui s'est produit depuis le jeudi à 18~h~00~; les deux autres ne copient que certains types d'événements,
et manquent alors tous les événements des autres types, quel que soit le moment de la semaine où ils se sont produits.
\end{solution}

\begin{solution}{exo:pl:trade-capture-and-booking:3}
Le LEI de la société, suivi d'une valeur unique pour cette société~: \texttt{5493001KJTIIGC8Y1R12} et \texttt{000\dots0007}, complété à 52 caractères,
ne contenant que des lettres majuscules et des chiffres uniquement.
\end{solution}

\begin{solution}{exo:pl:trade-capture-and-booking:4}
Sans cela, l'allocation peut favoriser un fonds après coup --- les meilleures exécutions, le lot irrégulier sur un marché haussier --- ce qui est inéquitable envers les
autres clients et, pour un gérant, constitue un conflit d'intérêts que les régulateurs recherchent. Une règle fixée à l'avance et appliquée mécaniquement supprime ce choix.
\end{solution}

\begin{solution}{exo:pl:trade-capture-and-booking:5}
Toutes les novations et toutes les résiliations partielles~: il ne recevait des copies que pour les nouvelles transactions et les modifications. Il a besoin de tous les événements qui modifient
ce qu'il confirme --- en pratique, le flux complet, filtré par lui-même plutôt que par l'émetteur.
\end{solution}

\begin{solution}{exo:pl:trade-capture-and-booking:6}
Comme un événement (un transfert de livre) avec son heure et sa raison, et non en modifiant le champ du livre~: le risque et le P\&L de la veille ont été déclarés dans
l'ancien livre et doivent rester reproductibles, et le transfert lui-même constitue une information pour le product control.
\end{solution}

\begin{solution}{exo:pl:trade-capture-and-booking:7}
Deux copies~: ajouter une copie à midi permet de capturer les événements du matin, il ne reste alors que ceux survenus entre midi et 15~h~00. Des copies plus fréquentes réduisent
la fenêtre~; elles ne la suppriment jamais.
\end{solution}

\begin{solution}{exo:pl:trade-capture-and-booking:8}
Lire à la demande depuis le système de trading n'évite les copies périmées que si ce système détient tous les événements et si chaque système lit la même
version au même instant~; en pratique, chacun lit à un moment différent, conserve ce qu'il a lu et manque ce que le système de trading ne gère pas
(novations, résiliations traitées ailleurs). Un accord suppose un seul flux d'événements et des versions datées d'un instant déclaré.
\end{solution}

\begin{solution}{pb:pl:trade-capture-and-booking:1}
\begin{enumerate}
\item Identifiants, contreparties, livre, dates, statut, version, ainsi que l'instrument, la quantité et le prix~; le format de la bibliothèque permet à chaque système de valoriser et
de gérer le risque de la transaction à partir du même \omterm{def:pl:trade-capture-and-booking:capture}{enregistrement}.
\item Le LEI de l'entité qui génère la transaction, suivi d'une valeur unique pour celle-ci, au plus 52 caractères, A--Z et 0--9.
\item Formats des identifiants, contreparties distinctes, règlement non antérieur à la date de transaction, quantité positive, type d'instrument.
\item Le livre dans lequel une transaction est enregistrée, déduit de son type d'instrument et du desk.
\item Pour que chaque version reste reproductible et que chaque système puisse reconstruire la même.
\item 69 modifications, 26 novations et 41 résiliations partielles parmi 200 swaps.
\item Modifiée le lundi, novée le mercredi, réduite de 40~\% le vendredi~; le risque détenait la version~3, la confirmation la version~2 et le règlement la version~4.
\item Risque~: les événements survenus depuis le dernier soir~; confirmation~: novations et résiliations~; règlement~: modifications et novations.
\item 25, 60 et 71 transactions.
\item 109 sur 200~; aucune.
\item Dix ordres enfants prenant le cours vendeur toutes les 20~secondes, 57 exécutions, prix moyen de 100.0783~\$.
\item 4\,950, 2\,970 et 1\,980 actions.
\item Plus grand reste 4\,900/3\,000/2\,000, reste 0~; lots irréguliers aux plus grands 5\,000/3\,000/1\,900, reste 0~; arrondi de chacun
5\,000/3\,000/2\,000, reste $-100$.
\item 50, 80 et 50 actions.
\item Le prix moyen du bloc, 100.0783~\$.
\item \textbf{Résultat nommé.} Après une semaine d'événements de cycle de vie, la distribution par copies laisse 25, 60 et 71 transactions erronées dans le risque, la confirmation et
le règlement et 109 swaps sur 200 pour lesquels ils divergent~; la distribution par événements, aucune. Le bloc de 9\,900 actions se répartit avec des restes de 0, 0
et $-100$ actions selon les règles du plus grand reste, des lots irréguliers aux plus grands et de l'arrondi de chacun, et des écarts à la cible d'au plus 50, 80 et 50 actions.
\item Plus grand reste~: il alloue la totalité du bloc et maintient chaque fonds au plus près de sa cible.
\item Qu'il consomme le flux d'événements et qu'il sache dire quelle version il détient, et à quel instant.
\item La part des transactions qui passent de l'exécution au règlement sans intervention manuelle.
\item Le résultat de ses événements.
\end{enumerate}
\end{solution}

\begin{solution}{iq:pl:trade-capture-and-booking:1}
Reconstruire la transaction à partir de ses événements à l'instant présent et comparer la version de chaque système avec elle~; remplacer ensuite les copies par des abonnements au flux d'événements,
chaque système indiquant la version qu'il détient.
\iqlookfor{Event sourcing et versions.}
\end{solution}

\begin{solution}{iq:pl:trade-capture-and-booking:2}
Un modèle canonique~: un en-tête d'identifiants, de contreparties, de livre, de dates, de statut et de version, et l'économie décrite par l'instrument de la bibliothèque de valorisation, avec
validation~; les systèmes conservent leurs propres vues dérivées, mais jamais leur propre copie des termes.
\iqlookfor{Un modèle unique, validé, partagé.}
\end{solution}

\begin{solution}{iq:pl:trade-capture-and-booking:3}
Au prix moyen du bloc, selon une règle fixée avant la transaction (au prorata des cibles, en lots entiers, les restes étant attribués par plus grande partie fractionnaire), le
bloc étant entièrement alloué, la règle et le résultat étant consignés.
\iqlookfor{Prix moyen, règle prédéfinie, aucun reste.}
\end{solution}

\begin{solution}{iq:pl:trade-capture-and-booking:4}
Le remplacement d'une contrepartie par une autre, avec le consentement de l'autre contrepartie~; la transaction se poursuit avec une nouvelle contrepartie, de sorte que l'exposition au crédit,
les confirmations, les instructions de règlement et la déclaration changent, alors que l'économie peut ne pas changer.
\iqlookfor{Changement de contrepartie comme événement aux effets en aval.}
\end{solution}

\begin{solution}{iq:pl:trade-capture-and-booking:5}
Des adaptateurs depuis chaque source d'exécution vers un modèle canonique~; des UTI et une validation à la saisie~; un \omterm{def:pl:trade-capture-and-booking:capture}{enregistrement} selon des règles~; des allocations selon des règles prédéfinies~;
des événements de cycle de vie dans un magasin en event sourcing~; chaque système aval abonné au flux d'événements, avec des versions datées~; un taux de \omterm{def:pl:trade-capture-and-booking:stp}{traitement de bout en bout} et
des ruptures surveillées.
\iqlookfor{Capturer une fois, événements partout.}
\end{solution}
